% =====================================================================
\begin{thebibliography}{99}
\addcontentsline{toc}{chapter}{Bibliografía}
\markboth{Bibliografía}{}
\small

\subsection*{Exámenes, hojas de problemas y trabajos de curso citados}

\bibitem{mit15053mid}
J.~Orlin y E.~Nasrabadi.
\newblock 15.053: Optimization Methods in Management Science. \emph{Midterm Exam~2: Practice Problem Set} y sus soluciones.
\newblock Massachusetts Institute of Technology, MIT OpenCourseWare, primavera de 2013. Problemas~2 («IP Formulation: Lockboxes») y~4 (restricción «al menos una de dos» y costes no lineales).
\newblock \url{https://ocw.mit.edu/courses/15-053-optimization-methods-in-management-science-spring-2013/pages/study-materials/}

\bibitem{mit15053ps}
J.~Orlin y E.~Nasrabadi.
\newblock 15.053: Optimization Methods in Management Science. Hojas de problemas~4, 5 y~6.
\newblock Massachusetts Institute of Technology, primavera de 2013.
\newblock \url{https://ocw.mit.edu/courses/15-053-optimization-methods-in-management-science-spring-2013/pages/assignments/}

\bibitem{cam4f13}
C.~E. Rasmussen.
\newblock 4F13: Probabilistic Machine Learning. Trabajo de curso~1, «Gaussian processes».
\newblock University of Cambridge, Michaelmas de 2024.
\newblock \url{https://mlg.eng.cam.ac.uk/teaching/4f13/2425/}

\bibitem{ee363final}
S.~Boyd.
\newblock EE363: Linear Dynamical Systems. Examen final (24 horas), invierno de 2006.
\newblock Stanford University. Problema~8, «LQR with linear state cost».
\newblock \url{https://ee363.stanford.edu/archive/notes/final_06.pdf}

\bibitem{ee363}
Stanford University.
\newblock EE363: Linear Dynamical Systems.
\newblock Primavera de 2026. Clases~16 «Linear quadratic regulator: Discrete-time finite horizon», 17 «LQR via Lagrange multipliers» y~18 «Infinite horizon linear quadratic regulator».
\newblock \url{https://ee363.stanford.edu/lectures.html}

\bibitem{mit18s096ps2}
A.~Edelman y S.~G. Johnson.
\newblock 18.S096: Matrix Calculus for Machine Learning and Beyond. Hoja de problemas~2.
\newblock Massachusetts Institute of Technology, MIT OpenCourseWare, IAP (enero) de 2023. Problema~1, gradiente de $f(p)=\bigl(c\T A(p)^{-1}b\bigr)^{2}$ con $A(p)$ tridiagonal mediante una resolución adjunta.
\newblock \url{https://ocw.mit.edu/courses/18-s096-matrix-calculus-for-machine-learning-and-beyond-january-iap-2023/}

\bibitem{mit1805final}
J.~Orloff y J.~Bloom.
\newblock 18.05: Introduction to Probability and Statistics. Examen final.
\newblock Massachusetts Institute of Technology, MIT OpenCourseWare, primavera de 2022. Problemas~9 (contraste $t$ e intervalo de confianza), 11 (potencia) y~13 (intervalo \emph{bootstrap} percentil).
\newblock \url{https://ocw.mit.edu/courses/18-05-introduction-to-probability-and-statistics-spring-2022/}

\bibitem{mit15401}
A.~Lo.
\newblock 15.401: Finance Theory I.
\newblock Massachusetts Institute of Technology, MIT OpenCourseWare, otoño de 2008. Sesiones «Present Value Relations» y «Capital Budgeting».
\newblock \url{https://ocw.mit.edu/courses/15-401-finance-theory-i-fall-2008/}

\subsection*{Libros y artículos}

\bibitem{bt}
D.~Bertsimas y J.~N. Tsitsiklis.
\newblock \emph{Introduction to Linear Optimization}.
\newblock Athena Scientific, 1997. Texto del curso 6.251J del MIT.

\bibitem{bv}
S.~Boyd y L.~Vandenberghe.
\newblock \emph{Convex Optimization}.
\newblock Cambridge University Press, 2004. Texto del curso EE364a de Stanford.
\newblock \url{https://web.stanford.edu/~boyd/cvxbook/}

\bibitem{kw}
M.~J. Kochenderfer y T.~A. Wheeler.
\newblock \emph{Algorithms for Optimization}.
\newblock MIT Press, 2019. Texto del curso AA222/CS361 de Stanford.
\newblock \url{https://algorithmsbook.com/optimization/}

\bibitem{jones}
D.~R. Jones, M.~Schonlau y W.~J. Welch.
\newblock Efficient global optimization of expensive black-box functions.
\newblock \emph{Journal of Global Optimization}, 13:455--492, 1998.

\bibitem{rw}
C.~E. Rasmussen y C.~K.~I. Williams.
\newblock \emph{Gaussian Processes for Machine Learning}.
\newblock MIT Press, 2006.
\newblock \url{https://gaussianprocess.org/gpml/}

\bibitem{garnett}
R.~Garnett.
\newblock \emph{Bayesian Optimization}.
\newblock Cambridge University Press, 2023.
\newblock \url{https://bayesoptbook.com/}

\bibitem{rmd}
J.~B. Rawlings, D.~Q. Mayne y M.~M. Diehl.
\newblock \emph{Model Predictive Control: Theory, Computation, and Design}, 2.ª ed.
\newblock Nob Hill Publishing, 2017. Caps.~1 y~2.
\newblock \url{https://sites.engineering.ucsb.edu/~jbraw/mpc/}

\bibitem{sb}
R.~S. Sutton y A.~G. Barto.
\newblock \emph{Reinforcement Learning: An Introduction}, 2.ª ed.
\newblock MIT Press, 2018.
\newblock \url{http://incompleteideas.net/book/the-book-2nd.html}

\bibitem{bs}
M.~P. Bendsøe y O.~Sigmund.
\newblock \emph{Topology Optimization: Theory, Methods, and Applications}, 2.ª ed.
\newblock Springer, Berlín, 2004. Cap.~1 «Topology optimization by distribution of isotropic material».

\bibitem{andreassen}
E.~Andreassen, A.~Clausen, M.~Schevenels, B.~S. Lazarov y O.~Sigmund.
\newblock Efficient topology optimization in MATLAB using 88 lines of code.
\newblock \emph{Structural and Multidisciplinary Optimization}, 43(1):1--16, 2011.

\bibitem{garland}
M.~Garland y P.~S. Heckbert.
\newblock Surface simplification using quadric error metrics.
\newblock En \emph{Proceedings of SIGGRAPH 97}, ACM, 1997.

\bibitem{bp}
R.~E. Barlow y F.~Proschan.
\newblock \emph{Mathematical Theory of Reliability}.
\newblock SIAM, 1996 (original: Wiley, 1965). Cap.~4, «Optimum Maintenance Policies».

\bibitem{montgomery}
D.~C. Montgomery.
\newblock \emph{Design and Analysis of Experiments}, 9.ª ed.
\newblock Wiley, 2017. Cap.~2 «Simple Comparative Experiments».

\subsection*{Cursos de ingeniería de sistemas y de sistemas con aprendizaje automático}

\bibitem{mit16842}
O.~de Weck.
\newblock 16.842: Fundamentals of Systems Engineering.
\newblock Massachusetts Institute of Technology, MIT OpenCourseWare, otoño de 2015.
\newblock \url{https://ocw.mit.edu/courses/16-842-fundamentals-of-systems-engineering-fall-2015/}

\bibitem{cs329s}
C.~Huyen y equipo docente de CS329S.
\newblock CS329S: Machine Learning Systems Design.
\newblock Stanford University, invierno de 2022.
\newblock \url{https://stanford-cs329s.github.io/syllabus.html}

\bibitem{huyen}
C.~Huyen.
\newblock \emph{Designing Machine Learning Systems}.
\newblock O'Reilly, 2022.

\bibitem{sculley}
D.~Sculley, G.~Holt, D.~Golovin, E.~Davydov, T.~Phillips, D.~Ebner, V.~Chaudhary, M.~Young, J.-F. Crespo y D.~Dennison.
\newblock Hidden technical debt in machine learning systems.
\newblock En \emph{Advances in Neural Information Processing Systems 28 (NeurIPS)}, 2015.

\subsection*{Apuntes del curso}

\bibitem{apuntes}
\emph{Anatomía de un LLM}. Apuntes de los módulos~24 a~28 (fase~VI) y de los módulos de las fases anteriores citados.
\newblock Edición de octubre de 2026.

\end{thebibliography}
